\documentclass[11pt,letterpaper]{article}
\usepackage[margin=0.86in,headheight=16pt]{geometry}
\usepackage[T1]{fontenc}
\usepackage{lmodern,microtype}
\usepackage{amsmath,amssymb,mathtools,bm}
\usepackage{graphicx,booktabs,array,siunitx}
\usepackage[dvipsnames]{xcolor}
\usepackage{fancyhdr,titlesec,caption,hyperref,placeins}
\usepackage{enumitem}
\usepackage{tikz}
\usetikzlibrary{arrows.meta,positioning,calc}
\hypersetup{colorlinks=true,linkcolor=MidnightBlue,urlcolor=MidnightBlue,citecolor=MidnightBlue}
\definecolor{deep}{RGB}{29,48,69}
\definecolor{accent}{RGB}{28,101,150}
\definecolor{muted}{RGB}{87,101,116}
\definecolor{pale}{RGB}{239,246,249}
\titleformat{\section}{\Large\bfseries\color{deep}}{\thesection}{0.7em}{}
\titleformat{\subsection}{\large\bfseries\color{accent}}{\thesubsection}{0.7em}{}
\titleformat{\subsubsection}{\normalsize\bfseries\color{deep}}{\thesubsubsection}{0.7em}{}
\fancyhf{}\pagestyle{fancy}
\fancyhead[L]{\small From atomic orbitals to quantum bonds}
\fancyhead[R]{\small Graduate-level notes}
\fancyfoot[C]{\thepage}
\setlength{\parskip}{4pt}\setlength{\parindent}{0pt}
\setlist[itemize]{leftmargin=1.3em,topsep=2pt,itemsep=2pt}
\newcommand{\ket}[1]{\lvert #1\rangle}
\newcommand{\bra}[1]{\langle #1\rvert}
\newcommand{\braket}[2]{\langle #1\mid #2\rangle}
\newcommand{\mel}[3]{\langle #1\lvert #2\rvert #3\rangle}
\newcommand{\ha}{\,\mathrm{Ha}}
\newcommand{\e}{\mathrm e}
\newcommand{\highlight}[1]{\par\smallskip\noindent\colorbox{pale}{\parbox{0.956\linewidth}{\strut #1\strut}}\smallskip\par}
\newcommand{\plot}[3]{\begin{figure}[htbp]\centering\includegraphics[width=#2\linewidth]{figures/#1.png}\caption{#3}\label{fig:#1}\end{figure}}
\newcommand{\rA}{r_A}\newcommand{\rB}{r_B}
\renewcommand{\arraystretch}{1.14}
\begin{document}
\begin{center}
{\fontsize{22}{27}\selectfont\bfseries From Atomic Orbitals to Quantum Bonds}\par\vspace{6pt}
{\Large Heitler--London, tunneling, exchange, and polarized $\sigma$/$\pi$ bonds}\par\vspace{8pt}
{\small A worked theoretical and numerical treatment using soluble one-dimensional models and hydrogenic three-dimensional orbitals}\par\vspace{6pt}
{\small Revised notes \textbullet\ November 2026}
\end{center}
\vspace{8pt}
\begin{abstract}
We start from neutral $\mathrm H_2$ and the Heitler--London spin-coupled electron pair. Two-site, hydrogenic, exactly solvable double-delta, and Gaussian variational models establish the difference between geometric overlap, Hamiltonian hopping, and direct and virtual charge-transfer exchange. The exact one-electron identity $H_{ab}=\langle a|h_B+V_A|b\rangle$ is used only where its eigenstate hypothesis holds. For carbon bonding, we test directed $sp^3$--$sp^3$ and transverse $p$--$p$ orbital ingredients against a full-valence semiempirical electronic-structure calculation for ethane and ethylene. We then construct the corresponding spin-coupled covalent and ionic valence-bond configurations, treating carbon--oxygen polarity in methanol and carbon--hydrogen bonds in methane. The one-electron calculation calibrates bond topology and qualitative polarization, not Heitler--London pair energies. We explicitly distinguish orbital overlap, localized orbital atomic populations, resonance weights, and actual bond energies. Long analytic integrals, numerical recipes, and model limitations appear in appendices.
\end{abstract}
\tableofcontents
\vspace{0.6em}
\highlight{\textbf{Interpretation rule.} An overlap integral measures geometry of a chosen orbital basis; an off-diagonal Hamiltonian element measures coupling; an orbital-energy splitting is not a bond dissociation energy; and the Coulomb interaction does not act directly on spin. The spin dependence of a bond follows from fermionic antisymmetry and electronic correlations.}

\section{Neutral $\mathrm H_2$: the Heitler--London starting point}
\subsection{A genuine two-electron Hamiltonian}
Fix two protons $A$ and $B$ separated by $R$ (Born--Oppenheimer approximation), and use atomic units $\hbar=m_e=e=a_0=1$. The \emph{physical} nonrelativistic Hamiltonian is
\begin{equation}\label{eq:physicalH2}
\hat H=\sum_{i=1}^{2}\left[-\tfrac12\nabla_i^2-\frac1{r_{iA}}-\frac1{r_{iB}}\right]+\frac1{r_{12}}+\frac1R.
\end{equation}
It commutes with the operator $P_{12}$ exchanging \emph{both} electron coordinates and spins. There is no separate attraction that belongs to ``electron 1'' or ``electron 2.'' 

Let $a(\mathbf r)=\pi^{-1/2}\e^{-r_A}$ and $b(\mathbf r)=\pi^{-1/2}\e^{-r_B}$ be normalized hydrogenic $1s$ orbitals. Their overlap is $S(R)=\braket a b$; it is nonzero at finite $R$. The two \emph{labeled spatial assignments} are
\begin{equation}
A(1,2)=a(1)b(2),\qquad B(1,2)=b(1)a(2),\qquad \braket A B=S^2.
\end{equation}
Neither assignment by itself is a permitted two-electron fermionic state. To impose antisymmetry, introduce the spin singlet $\chi_s=[\alpha(1)\beta(2)-\beta(1)\alpha(2)]/\sqrt2$ and the three symmetric triplets $\chi_{t,m}$.

\subsection{Spin dictates spatial symmetry, not the energetic ordering}
The spatial combinations and full fermionic wavefunctions are
\begin{equation}\label{eq:HLstates}
\psi_{\pm}(1,2)=\frac{A(1,2)\pm B(1,2)}{\sqrt{2(1\pm S^2)}},\qquad
\boxed{\Psi_s=\psi_+\chi_s,\quad\Psi_{t,m}=\psi_-\chi_{t,m}.}
\end{equation}
The singlet is antisymmetric in spin and symmetric in space; the triplet is symmetric in spin and antisymmetric in space. Consequently $\psi_-(\mathbf r,\mathbf r)=0$, an exchange node imposed by Pauli exclusion. The symmetric wavefunction need not vanish at coincident positions. No spin-dependent force was added.

\plot{01_h2_atomic_orbital_axis}{0.96}{Hydrogenic $1s$ amplitudes along the bond axis and normalized symmetric/antisymmetric one-electron combinations. These are 3D orbitals shown as \emph{axial cuts}, not 1D normalized wavefunctions.}
\plot{02_h2_exchange_nodes}{0.97}{Axial sections of the two-electron \emph{spatial} Heitler--London combinations, $\psi_\pm(z_1,z_2)$. Red/blue indicate opposite signs, not probabilities. The triplet's exchange node lies on $z_1=z_2$; the full nodal condition is $\mathbf r_1=\mathbf r_2$.}

\subsection{The direct and exchange matrix elements}
Define $D=\mel A{\hat H}A$ and $X=\mel A{\hat H}B$. Electron-exchange symmetry forces the projected matrices to have the form
\begin{equation}\label{eq:HL-generalized}
\mathsf H=\begin{pmatrix}D&X\\X&D\end{pmatrix},\quad
\mathsf S=\begin{pmatrix}1&S^2\\S^2&1\end{pmatrix},\quad
\mathsf H\mathbf c=E\mathsf S\mathbf c.
\end{equation}
The resulting \emph{variational} energies in the restricted covalent space are
\begin{equation}\label{eq:HLenergies}
\boxed{E_s=\frac{D+X}{1+S^2},\qquad E_t=\frac{D-X}{1-S^2},\qquad
E_t-E_s=\frac{2(DS^2-X)}{1-S^4}.}
\end{equation}
At $R\to\infty$, $S\to0$ and both energies tend to $-1\ha$, the energy of two isolated H atoms. A true molecule is bound only when its \emph{total} minimum falls below that separated-fragment energy. The sign of $X$ \emph{alone} does not determine which spin state is favored; the overlap denominator and direct energy also enter.

There is a terminological trap: $X=\mel{ab}{H}{ba}$ is an \emph{exchange matrix element between two configurations}, while $K_{ee}=\iint a(1)b(1)r_{12}^{-1}a(2)b(2)\,d1d2$ is specifically the positive two-electron Coulomb exchange integral. The former additionally contains one-electron Hamiltonian and nuclear terms; they are not interchangeable.

\subsection{The covalent Heitler--London state versus a doubly occupied MO}
Define a normalized bonding orbital $g=(a+b)/\sqrt{2(1+S)}$. Restricted closed-shell HF uses $g(1)g(2)\chi_s$. Its spatial product contains \emph{both} covalent terms $ab+ba$ and ionic terms $aa+bb$, with constrained coefficients. By contrast, Eq.~\eqref{eq:HLstates} contains only the covalent assignment pair. At dissociation, the latter naturally gives one electron at each nucleus, whereas a single closed-shell bonding MO retains undesired ionic weight. The distinction is between \emph{many-electron wavefunctions}, not merely two orbital drawing conventions.

\section{Overlap and hopping}
\subsection{Two position states: the bare minimum model}
For one electron, begin with \emph{orthonormal} discrete states $\{\ket L,\ket R\}$ representing sites. Choose two normalized, nonorthogonal \emph{atomic trial vectors}
\begin{equation}\label{eq:toy-epsilon}
\ket a=\frac{\ket L+\epsilon\ket R}{\sqrt{1+\epsilon^2}},\qquad
\ket b=\frac{\epsilon\ket L+\ket R}{\sqrt{1+\epsilon^2}},\qquad
S=\braket a b=\frac{2\epsilon}{1+\epsilon^2}.
\end{equation}
This $\epsilon$ measures leakage \emph{relative to the chosen orthogonal site basis}. It does not determine an energy. A spin-independent effective Hamiltonian with explicitly specified hopping is
\begin{equation}\label{eq:toyh}
H_{\rm site}=\begin{pmatrix}E_0&-t\\-t&E_0\end{pmatrix}_{L,R},\qquad
\ket{\pm}=\frac{\ket L\pm\ket R}{\sqrt2},\quad E_\pm=E_0\mp t.
\end{equation}
For this special symmetric parameterization, $(a\pm b)/\sqrt{2(1\pm S)}=(L\pm R)/\sqrt2$; the \emph{same} delocalized eigenvectors follow for any $\epsilon$. If $t=0$ the states remain degenerate regardless of $S$. Hence \emph{overlap does not cause splitting without a Hamiltonian coupling}.
\plot{04_discrete_overlap_and_hopping}{0.93}{Left: leakage $\epsilon$ and atomic overlap $S$ are geometrical properties of a basis. Right: the eigenvalue splitting of the orthogonal site Hamiltonian is set by the independent parameter $t$.}

\subsection{Why the off-diagonal element needs overlap correction}
In a general normalized, nonorthogonal pair $\{a,b\}$, write $H_{aa}=H_{bb}=\alpha$, $H_{ab}=\beta$, $\braket a b=S$. The generalized eigenproblem gives
\begin{equation}
E_\pm=\frac{\alpha\pm\beta}{1\pm S},\quad
\boxed{t_{\rm orth}=\frac{E_--E_+}{2}=\frac{\alpha S-\beta}{1-S^2}.}
\label{eq:orthhop}
\end{equation}
This is the hopping coefficient after symmetric (L\"owdin) orthogonalization. In particular $H=E_0 I$ gives $\beta=E_0S\neq0$ but $t_{\rm orth}=0$. The ordinary matrix element $\beta$ is therefore \emph{not} itself the observable tunneling amplitude. For unequal sites, one must retain both diagonal energies and solve the generalized eigenvalue problem; the equal-site splitting formula does not apply directly.

\section{$\mathrm H_2^+$: physical Coulomb integrals and the decomposition $H=H_B+V_A$}
Before returning to two electrons, it helps to derive $t(R)$ for one electron using a real Coulomb potential. Write
\begin{equation}\label{eq:h2plus-split}
\hat h=-\tfrac12\nabla^2-\frac1{r_A}-\frac1{r_B}
=\underbrace{\left(-\tfrac12\nabla^2-\frac1{r_B}\right)}_{\hat h_B}+\underbrace{\left(-\frac1{r_A}\right)}_{V_A}.
\end{equation}
\emph{Do not} write $h=h_A+h_B$: this would count $T=-\frac12\nabla^2$ twice. Since $h_Bb=-\frac12b$,
\begin{equation}\label{eq:habsplit}
\boxed{H_{ab}=\mel a{h_B+V_A}b=E_B\braket a b+\mel a{V_A}b.}
\end{equation}
The first contribution comes from acting with the exactly solvable right-atom Hamiltonian; the second is the new attraction from the left proton. It is valid only because $b$ is genuinely an eigenfunction of $h_B$.

For normalized $1s$ orbitals the required integrals can be done analytically (Appendix~\ref{app:1s}):
\begin{align}\label{eq:h2plus-integrals}
S(R)&=\int ab\,d^3r=\e^{-R}(1+R+R^2/3),\\
J_N(R)&=\int\frac{|a|^2}{r_B}\,d^3r=\frac{1-(1+R)\e^{-2R}}{R},\\
K_N(R)&=\int\frac{ab}{r_A}\,d^3r=(1+R)\e^{-R}.\nonumber
\end{align}
The diagonal and off-diagonal Hamiltonian elements are
\begin{equation}\label{eq:h2plus-matrix}
\alpha=-\tfrac12-J_N,\qquad\beta=-\tfrac12 S-K_N.
\end{equation}
Therefore, after accounting for overlap,
\begin{equation}\label{eq:h2plus-t}
\boxed{t(R)=\frac{\alpha S-\beta}{1-S^2}=\frac{K_N-J_NS}{1-S^2}.}
\end{equation}
The normalized MO energies are $E_\pm^{(e)}=(\alpha\pm\beta)/(1\pm S)$. The \emph{total} one-electron $\mathrm H_2^+$ energies include nuclear repulsion: $E_\pm^{\mathrm{total}}=E_\pm^{(e)}+1/R$. At $R=2a_0$, $S=0.586453$ and $t=0.196459\ha$, but $E_+^{\mathrm{total}}=-0.553771\ha$. The total binding at that fixed separation relative to $\mathrm H+\mathrm H^+$ is only $0.053771\ha$; it is not equal to $t$.
\plot{03_h2plus_overlap_energy}{0.97}{Hydrogenic fixed-$1s$ $\mathrm H_2^+$ variational model. The four positive quantities on the left play different roles; the energy plot on the right includes $1/R$ proton repulsion.}
\highlight{\textbf{Why this is useful later.} The identity $H_{ab}=E_BS+\langle a\lvert V_A\rvert b\rangle$ works for \emph{any} atomic eigenfunction $b$, including hydrogenic $2s$, $2p$, and an $sp^3$ combination within a degenerate $n=2$ manifold. It will fail as an equality with $E_BS$ for an arbitrary Gaussian trial function that is not an eigenstate of $H_B$.}
\FloatBarrier

\section{An exactly solvable one-dimensional molecule: two attractive delta wells}
\subsection{An isolated ``atom'' with an exact exponential orbital}
Use the physically localized one-dimensional Hamiltonian
\begin{equation}
\hat h_A=-\frac{\hbar^2}{2m}\frac{d^2}{dx^2}-g\delta(x+R/2),\quad
\hat h_B=-\frac{\hbar^2}{2m}\frac{d^2}{dx^2}-g\delta(x-R/2),\quad g>0.
\end{equation}
Each has an \emph{exact} bound eigenstate with $\kappa=mg/\hbar^2$:
\begin{equation}\label{eq:delta-atoms}
a(x)=\sqrt\kappa\,\e^{-\kappa|x+R/2|},\quad
b(x)=\sqrt\kappa\,\e^{-\kappa|x-R/2|},\quad
E_0=-\frac{\hbar^2\kappa^2}{2m}.
\end{equation}
The inverse-potential idea is rigorous here in the distributional sense: the derivative jump $\psi'(x_A^+)-\psi'(x_A^-)=-2\kappa\psi(x_A)$ produces the attractive delta function. The normalized atomic overlap is
\begin{equation}\label{eq:delta-overlap}
\boxed{S(R)=\int a(x)b(x)\,dx=(1+\kappa R)\e^{-\kappa R}.}
\end{equation}

\subsection{Full double well: exact finite-$R$ states}
The actual molecular Hamiltonian is $\hat h_R=T+V_A+V_B=h_B+V_A$, \emph{not} $h_A+h_B$. Its even and odd bound states are
\begin{equation}\label{eq:delta-exact}
\boxed{\psi_\pm^{\rm exact}(x;R)=N_\pm\Big(\e^{-q_\pm|x+R/2|}\pm\e^{-q_\pm|x-R/2|}\Big),}
\end{equation}
where $E_\pm=-\hbar^2q_\pm^2/(2m)$ and the jump conditions imply
\begin{equation}\label{eq:delta-transcendental}
\boxed{q_+=\kappa(1+\e^{-q_+R}),\qquad q_-=\kappa(1-\e^{-q_-R}).}
\end{equation}
The normalization is $N_\pm^2=q_\pm/[2\{1\pm(1+q_\pm R)\e^{-q_\pm R}\}]$. The odd \emph{bound} state exists only for $\kappa R>1$; otherwise the odd sector has no discrete normalizable eigenstate.
\plot{05_delta_exact_vs_lcao}{0.97}{Numerically solved exact double-delta eigenstates (solid) compared to normalized fixed-$\kappa$ LCAO combinations (dashed), at three separations. Here $\hbar=m=g=\kappa=1$. All plotted states exist because $R>1$.}
\plot{06_delta_decay_energies}{0.97}{The exact decay constants and energies. The two wells become independent as $R\to\infty$, while the odd bound level reaches the continuum at $R=1/\kappa$.}

\subsection{A completely explicit origin for hopping}
Use the right-well eigenvalue equation $h_Bb=E_0b$. Our preferred split immediately gives
\begin{align}
H_{ab}&=\mel a{h_B+V_A}b=E_0S-g\,a(-R/2)b(-R/2)\notag\\
&=\boxed{E_0S-g\kappa\e^{-\kappa R}},\label{eq:delta-hab}\\
H_{aa}&=\mel a{h_A+V_B}a=\boxed{E_0-g\kappa\e^{-2\kappa R}}.\label{eq:delta-haa}
\end{align}
The resulting \emph{projected} hopping coefficient is
\begin{equation}\label{eq:delta-t}
\boxed{t_{\rm proj}(R)=\frac{H_{aa}S-H_{ab}}{1-S^2}
=g\kappa\frac{\e^{-\kappa R}-S\e^{-2\kappa R}}{1-S^2}.}
\end{equation}
For $\kappa R\gg1$, $t_{\rm proj}\sim g\kappa\e^{-\kappa R}=(\hbar^2\kappa^2/m)\e^{-\kappa R}$. This agrees at leading order with half the exact even--odd splitting obtained from Eq.~\eqref{eq:delta-transcendental}. The \emph{finite-$R$ projected} splitting and the \emph{exact continuum} splitting are different quantities because the fixed atomic LCAO space does not allow $q_\pm$ to relax.

\section{A Gaussian basis as a genuine variational worked example}
\subsection{First fix the Hamiltonian; then choose an imperfect trial function}
Keep the \emph{same} attractive delta Hamiltonian $h_A=-\frac12\partial_x^2-g\delta(x-x_A)$, now in atomic units. Take a normalized Gaussian with width $\sigma$:
\begin{equation}\label{eq:gauss}
g_A(x;\sigma)=(2\pi\sigma^2)^{-1/4}\exp[-(x+R/2)^2/(4\sigma^2)],\quad
g_B(x;\sigma)=g_A(-x;\sigma).
\end{equation}
Each orbital is normalized \emph{independently of $R$}, and their overlap follows from elementary Gaussian integration:
\begin{equation}\label{eq:gauss-S}
\boxed{S_G(R,\sigma)=\braket{g_A}{g_B}=\exp[-R^2/(8\sigma^2)].}
\end{equation}
Unlike the exponential delta-well ground state, the Gaussian is \emph{not} an eigenfunction of the delta Hamiltonian. That is fine: the Rayleigh--Ritz principle applies to any normalized trial state.

For a \emph{single} delta well, $\langle T\rangle=1/(8\sigma^2)$ and $\langle-g\delta(x-x_A)\rangle=-g/(\sqrt{2\pi}\sigma)$, so
\begin{equation}\label{eq:gauss-variational}
\boxed{E_G(\sigma)=\frac1{8\sigma^2}-\frac{g}{\sqrt{2\pi}\sigma},\quad
\sigma_{\rm opt}=\frac{\sqrt{2\pi}}{4g},\quad E_G^{\min}=-\frac{g^2}{\pi}.}
\end{equation}
Compare the exact isolated-delta energy $E_{\rm exact}=-g^2/2$. The variational value is higher, as it must be: $-1/\pi\simeq-0.3183>-0.5$ at $g=1$.

\subsection{Gaussian matrix elements in the two-center Hamiltonian}
The kinetic cross integral is $\langle g_A\lvert T\rvert g_B\rangle=S_G(1/8\sigma^2-R^2/32\sigma^4)$. The two delta potentials give identical cross contributions. The full matrices are therefore
\begin{align}\label{eq:gauss-mat}
H_{AA}^{(G)}&=\frac1{8\sigma^2}-\frac{g}{\sqrt{2\pi}\sigma}\Big(1+\e^{-R^2/(2\sigma^2)}\Big),\\
H_{AB}^{(G)}&=S_G\Big(\frac1{8\sigma^2}-\frac{R^2}{32\sigma^4}\Big)
-\frac{2g}{\sqrt{2\pi}\sigma}\e^{-R^2/(4\sigma^2)}.\nonumber
\end{align}
Thus $E_\pm^{(G)}(R,\sigma)=[H_{AA}^{(G)}\pm H_{AB}^{(G)}]/[1\pm S_G]$, which can be minimized over $\sigma$ \emph{separately in each parity sector}. At $g=1$, $R=2.5$, the optimized Gaussian energies are $E_+^{(G)}\simeq-0.42107$ and $E_-^{(G)}\simeq-0.19597$, both above the corresponding exact bound-state levels.
\plot{07_gaussian_variational}{0.97}{Left: the Gaussian is a legitimate variational approximation to an isolated delta well, not an exact eigenstate. Right: independently width-optimized even and odd Gaussian LCAO states compared with the exact double-delta bound states at $R=2.5$.}

\subsection{Precisely when can we write $H_{ab}=E_B S+\langle a\lvert V_A\rvert b\rangle$?}
The operator decomposition $H=H_B+V_A$ is \emph{always} true for the two-well Hamiltonian. But the next step requires $H_B\ket b=E_B\ket b$. For a Gaussian trial state $\ket{g_B}$,
\begin{equation}
\mel{g_A}{H}{g_B}=E_BS_G+\mel{g_A}{V_A}{g_B}
+\underbrace{\mel{g_A}{H_B-E_B}{g_B}}_{\text{nonzero variational residual}}.
\end{equation}
If we silently discard the last term, we incorrectly treat a Gaussian trial orbital as an atomic eigenstate. This is the essential distinction between an \emph{exactly soluble reference Hamiltonian} and a \emph{variational basis}.
\FloatBarrier
\section{Exchange interactions: what actually favors a singlet?}
\subsection{Two distinct mechanisms with opposite tendencies}
The phrase \emph{exchange interaction} is overloaded. Given two \emph{orthogonal} localized orbitals $a,b$, the repulsive Coulomb operator alone contributes
\begin{equation}
J_{ee}=\iint\frac{|a(1)|^2|b(2)|^2}{r_{12}}\,d1d2,\qquad
K_{ee}=\iint\frac{a^*(1)b(1)b^*(2)a(2)}{r_{12}}\,d1d2\geq0.
\end{equation}
Its singlet and triplet expectations in the symmetric/antisymmetric \emph{covalent spatial} states are $J_{ee}+K_{ee}$ and $J_{ee}-K_{ee}$. Thus \emph{direct Coulomb exchange} by itself favors a triplet, with $E_t-E_s=-2K_{ee}$. This is related to Hund's-rule exchange. It does \emph{not} mean that the physical $\mathrm H_2$ triplet is the ground state; electron--nucleus attraction, kinetic delocalization, orbital overlap, and virtual charge transfer all participate in the full problem.

For spatially overlapping $a,b$, one must retain the $1\pm S^2$ normalization in Eq.~\eqref{eq:HLenergies}. The difference between an exchange \emph{integral} and a singlet--triplet \emph{energy splitting} is fundamental.

\subsection{The Hubbard dimer: legitimate degenerate perturbation theory with fermions}
Pass to an \emph{orthonormal} localized site basis, and introduce two electrons through fermionic creation operators $c^\dagger_{A\sigma},c^\dagger_{B\sigma}$. The exact finite-dimensional Hamiltonian is
\begin{equation}\label{eq:hubbard}
\hat H=-t\sum_{\sigma=\uparrow,\downarrow}(c^\dagger_{A\sigma}c_{B\sigma}+c^\dagger_{B\sigma}c_{A\sigma})
+U\sum_{j=A,B}n_{j\uparrow}n_{j\downarrow},\qquad U>0.
\end{equation}
The covalent sector (one electron per site) contains one properly antisymmetrized singlet and three triplets. The singlet and symmetric ionic state can be written in first-quantized spatial notation, each with singlet spin:
\begin{equation}
\ket C_s=\frac{\ket{AB}+\ket{BA}}{\sqrt2},\qquad
\ket I_s=\frac{\ket{AA}+\ket{BB}}{\sqrt2},\qquad
\ket C_t=\frac{\ket{AB}-\ket{BA}}{\sqrt2}\quad(\text{triplet spin}).
\end{equation}
Let $T$ denote the hopping term. Its spatial action is $T\ket{AB}=-t(\ket{AA}+\ket{BB})$ and the same for $\ket{BA}$. Consequently,
\begin{equation}\label{eq:hubbard-2t}
\boxed{T\ket C_s=-2t\ket I_s,\qquad T\ket C_t=0.}
\end{equation}
There is no singlet--triplet \emph{matrix element}: their spin wavefunctions are orthogonal, and the spin-independent $T$ cannot change total spin. The singlet instead couples to a \emph{higher-energy state in its own spin sector}; the triplet cannot make the corresponding hop.

Choose $H_0=U\sum_jn_{j\uparrow}n_{j\downarrow}$ and $V=T$. Now the zero-energy covalent singlet and three triplets are \emph{exactly degenerate eigenstates of an exchange-symmetric $H_0$}. Moreover, $PVP=0$ in that four-state covalent manifold. Hence first-order degenerate perturbation theory yields no splitting. At second order,
\begin{equation}\label{eq:hubbard-second}
E_s^{(2)}=\frac{|\mel{I_s}{T}{C_s}|^2}{0-U}=-\frac{4t^2}{U},\quad
E_t^{(2)}=0,\quad
\boxed{J_{\mathrm{AF}}\equiv E_t-E_s\simeq\frac{4t^2}{U}>0.}
\end{equation}
The exact singlet matrix in the basis $\{\ket C_s,\ket I_s\}$ is
\begin{equation}\label{eq:hubbard-exact}
H_s=\begin{pmatrix}0&-2t\\-2t&U\end{pmatrix},\qquad
E_s^-=\frac{U-\sqrt{U^2+16t^2}}{2},\quad E_t=0.
\end{equation}
The other singlets have energies $(U+\sqrt{U^2+16t^2})/2$ and $U$; the latter is the antisymmetric ionic configuration. In the correlated ground state, the ionic probability is $p_I=\tfrac12[1-U/\sqrt{U^2+16t^2}]$.
\plot{08_hubbard_exchange}{0.96}{The exactly solvable two-electron Hubbard dimer. The energy gain begins at order $t^2/U$, and the associated ionic admixture grows as charge-transfer configurations become accessible.}

\subsection{Effective spin Hamiltonian and the sign of exchange}
Within the singly occupied subspace, $\bm S_A\cdot\bm S_B=-3/4$ in the singlet and $+1/4$ in the triplets. Up to a common energy shift,
\begin{equation}\label{eq:heisenberg}
\boxed{H_{\rm spin}=J_{\rm eff}\left(\bm S_A\cdot\bm S_B-\tfrac14\right),\quad J_{\rm eff}=\frac{4t^2}{U}\quad\text{(Hubbard, second order)}.}
\end{equation}
The repulsion-only direct-exchange contribution would instead yield $J_{\rm direct}\sim-2K_{ee}$ in a simple orthogonal-orbital approximation. A broader model can contain both, schematically $J_{\rm eff}\approx4t^2/U-2K_{ee}+\cdots$; the ellipsis includes intersite Coulomb, correlated hopping, orbital relaxation, and other terms. These contributions must be derived from \emph{one consistently projected Hamiltonian} to avoid double counting.

\highlight{\textbf{Connection to $\mathrm H_2$.} The Heitler--London splitting in Eq.~\eqref{eq:HLenergies} is obtained by diagonalizing a covalent two-configuration variational space. The Hubbard superexchange $4t^2/U$ comes from \emph{second-order} excursions into ionic states that the simplest Heitler--London ansatz excludes. Both are manifestations of fermionic exchange symmetry, but they are not the same approximation or the same integral.}
\FloatBarrier

\section{From atomic pair functions to chemically useful valence bonds}
\subsection{What went wrong with the original higher-orbital examples}
Our earlier two-center construction treated $Z_{\rm eff}$ as if it were a physical point charge at \emph{every} radius: $h=T-Z_A/r_A-Z_B/r_B$. For C and O, the parameters $Z_{\rm eff}$ describe screening by \emph{other electrons}; they are neither bare nuclear charges nor adequate all-radius effective potentials. The resulting model left out the carbon and oxygen occupied cores, oxygen lone pairs, and the remaining valence electrons. At the selected C--O distance it predicted a $98.8\%$ oxygen weight in a one-electron state; for C--H it perversely put $85.7\%$ on H. These are not sensible predictions of bond polarity.

A second failure was to compare hydrogenic $2s$/$2p$ functions with radial nodes at separations where cancellation can overwhelm the head-on overlap of the $sp^3$ hybrids. An \emph{isolated one-electron two-orbital model} cannot be trusted to represent a C--C $\sigma$ bond in a molecule that has six further C--H bonds, or a C--O bond in the presence of O lone pairs.

Our replacement has two deliberately separate levels:
\begin{enumerate}
\item A \textbf{full-valence semiempirical electronic structure} calculation supplies an interpretable, geometry-dependent occupied space and tests the $\sigma/\pi$ pattern, bond indices, and polarization. It is \emph{not} a valence-bond calculation and does not yield trustworthy bond energies.
\item A \textbf{spin-coupled valence-bond (VB) construction} uses the atomic/hybrid orbitals as \emph{ingredients} of covalent and ionic two-electron configurations. Fermionic antisymmetry and the Coulomb interaction determine the pair energies, as in Heitler--London $\mathrm H_2$. We keep the necessary many-electron integrals symbolic unless they have genuinely been evaluated.
\end{enumerate}
This division preserves the physical mechanism without inventing quantitative VB energies from one-electron orbital eigenvalues.

\subsection{Nodeless, normalized valence shapes with angular structure}
For carbon-like and oxygen-like valence functions use the following \emph{Slater-type orbitals} (STOs), with an adjustable positive radial exponent $\zeta$. Unlike the exact hydrogenic $2s$ state, the nodeless radial $2s$ function is a compact shape approximation to a core-orthogonalized valence orbital:
\begin{equation}\label{eq:sto-valence}
\boxed{\begin{aligned}
s_\zeta(\mathbf r)&=\sqrt{\frac{\zeta^5}{3\pi}}\,r\e^{-\zeta r},\\
p_{j,\zeta}(\mathbf r)&=\sqrt{\frac{\zeta^5}{\pi}}\,r_j\e^{-\zeta r},\quad j=x,y,z,\\
h_{sp^3,\hat n}&=\tfrac12\bigl(s_\zeta+\sqrt3\,p_{\hat n,\zeta}\bigr),\qquad
h_{sp^2,\hat n}=\tfrac1{\sqrt3}\bigl(s_\zeta+\sqrt2\,p_{\hat n,\zeta}\bigr).
\end{aligned}}
\end{equation}
The functions are orthonormal \emph{on one center}. The $p_{\hat n}$ orbital is the $p$ function directed along the chosen bond; for an opposed pair we take $\hat n_A=+\hat z$ and $\hat n_B=-\hat z$. A transverse $p_x$ or $p_y$ instead supplies a $\pi$-pair ingredient. For hydrogen use $1s_\zeta=(\zeta^3/\pi)^{1/2}\e^{-\zeta r}$. The exponents set orbital \emph{size} and are not literal point-charge values inserted into the potential.

The empirical parameters used for the reproducible valence calculation are
\begin{center}\small
\begin{tabular}{lrrrr}
\toprule
Atom & $\zeta$ ($a_0^{-1}$) & $H_{ss}$ (eV) & $H_{pp}$ (eV) & Valence electrons\\\midrule
H (1s) & 1.300 & $-13.60$ & --- & 1\\
C (2s, 2p) & 1.625 & $-21.40$ & $-11.40$ & 4\\
O (2s, 2p) & 2.275 & $-32.38$ & $-15.85$ & 6\\\bottomrule
\end{tabular}
\end{center}
The oxygen-to-carbon valence exponent ratio $\zeta_O/\zeta_C=2.275/1.625=1.400$ is a \emph{model radial-size ratio}, not a fitted C--O partial charge. It retains the originally requested unequal-size $1.4$ comparison without treating screened effective charges as a physical molecular Coulomb Hamiltonian.

\subsection{Full valence space as a calibration, not as an alternative to VB}
Use all H $1s$ and C/O $2s,2p$ valence functions for methane, ethane, ethylene, and methanol, at the stated molecular geometries. Compute every interatomic overlap $S_{\mu\nu}$ by three-dimensional quadrature; keep $S_{\mu\mu}=1$ and orthogonal $s/p$ functions on the same atom. In the extended H\"uckel approximation,
\begin{equation}\label{eq:eht}
H_{\mu\mu}=\varepsilon_\mu^{\rm atom},\qquad
\boxed{H_{\mu\nu}=K S_{\mu\nu}\frac{\varepsilon_\mu^{\rm atom}+\varepsilon_\nu^{\rm atom}}2\quad (\mu\ne\nu),\quad K=1.75,}
\end{equation}
followed by $HC=SC\varepsilon$. These are semiempirical valence-state parameters; the model includes neither self-consistent electron repulsion nor a repulsive nuclear/core total-energy functional. The supplied bond lengths are \emph{geometry inputs}, not predicted equilibrium lengths.

For completeness, a normalized occupied orbital rotation $C_{\rm loc}=C_{\rm occ}U$ with $U^\dagger U=I$ leaves the occupied density unchanged. Our localization criterion maximizes $\sum_{i,A}q_{Ai}^4$, where $q_{Ai}=\sum_{\mu\in A}|(S^{1/2}C_{\rm occ}U)_{\mu i}|^2$. This is a \emph{L\"owdin-population localization}, not the intrinsic-atomic-orbital (IAO) procedure of genuine IBO localization.

A useful orbital-bond index, in the same population convention, is
\begin{equation}\label{eq:wiberg}
W_{AB}=\sum_{\mu\in A,\nu\in B}\bigl(P_{\rm orth}\bigr)_{\mu\nu}^2,\quad
P_{\rm orth}=S^{1/2}(2C_{\rm occ}C_{\rm occ}^{T})S^{1/2}.
\end{equation}
The numbers measure connectivity within this Hamiltonian and population partition; $W=1$ or $2$ is \emph{not} a VB bond dissociation energy.

\subsection{Return immediately to the Heitler--London electron pair}
For \emph{each} candidate bond $k$, choose two atom-centered orbitals $a_k,b_k$ (orthonormalization between different bonds and Pauli antisymmetrization across pairs are still necessary in a molecule). The normalized \emph{covalent} spin-singlet pair is exactly analogous to $\mathrm H_2$:
\begin{equation}\label{eq:VBpair-general}
\boxed{\Phi_{k,\mathrm{cov}}^{S=0}=\frac{a_k(1)b_k(2)+b_k(1)a_k(2)}{\sqrt{2(1+S_k^2)}}
\,\frac{\alpha(1)\beta(2)-\beta(1)\alpha(2)}{\sqrt2},\qquad S_k=\braket{a_k}{b_k}.}
\end{equation}
The corresponding triplet has the \emph{antisymmetric} spatial numerator and $\sqrt{2(1-S_k^2)}$ denominator. A polarized or correlation-corrected bond also requires ionic singlet configurations $|aa\rangle\chi_s$ and $|bb\rangle\chi_s$, with coefficients determined from a \emph{two-electron} generalized eigenproblem (the three configurations are nonorthogonal when $S_k\ne0$).

For a \emph{many-bond} molecule, an approximate single-structure spin-paired VB state is $\mathcal A[\Phi_{\sigma,1}\Phi_{\sigma,2}\cdots\Phi_{\pi}\times\text{lone-pair functions}]$, where $\mathcal A$ antisymmetrizes \emph{all} electrons. Additional VB structures and independent orbital optimization may be required. The localized occupied orbitals from Eq.~\eqref{eq:eht} show \emph{where useful pair regions lie}; they do not automatically supply the correlated VB coefficients.

\subsection{Why we no longer apply $H_{ab}=E_bS+\mel a{V_A}b$ without checking the reference}
Our exact hydrogen and delta-well identity remains the cleanest demonstration of physical hopping:
\begin{equation}
H_{ab}=\mel a{h_B+V_A}b=E_bS+\mel a{V_A}b
\quad\text{\emph{if} }h_Bb=E_bb.
\end{equation}
A Slater-type \emph{trial} orbital in an empirical effective environment is generally \emph{not} an eigenstate of any specified bare $h_B=T-Z_B/r_B$. The exact decomposition instead contains the residual
\begin{equation}\label{eq:sto-residual}
\boxed{H_{ab}=E_bS+\mel a{V_A}b+\mel a{h_B-E_b}b.}
\end{equation}
For the higher-orbital examples we therefore \emph{do not} use the Coulomb split to fabricate $H_{ab}$. We evaluate $S$ spatially and specify the empirical Hamiltonian matrix elements independently through Eq.~\eqref{eq:eht}. This matters: orbital overlap is geometric, while hopping and exchange are Hamiltonian dependent.
\FloatBarrier

\section{Equivalent $sp^3$--$sp^3$ $\sigma$ pairs: shape model and ethane}
\subsection{An equal-$Z=1$ angular experiment}
To isolate angular symmetry from chemical complexity, take two identical, nodeless STO valence functions with $\zeta_A=\zeta_B=Z/2=0.5\,a_0^{-1}$, and oppositely directed $sp^3$ hybrids along the internuclear $z$ axis. Set $Z=1$ \emph{only as a common radial-scale label}; these nodeless orbitals are not the exact excited eigenstates of a hydrogenic $Z=1$ atom. The functions are
\begin{equation}\label{eq:sto-sigma}
\boxed{a_\sigma=\frac{N_s}{2}\e^{-\zeta r_A}(r_A+3z_A),\qquad
b_\sigma=\frac{N_s}{2}\e^{-\zeta r_B}(r_B-3z_B),\quad N_s=\sqrt{\frac{\zeta^5}{3\pi}}.}
\end{equation}
Expand $S_\sigma=\int a_\sigma b_\sigma\,d^3r$ into $ss$, $sp$, $ps$, and $pp$ overlap terms. The $p$ signs select \emph{facing} lobes; unlike the old radially nodal hydrogenic example, this shape model shows the expected head-on overlap at a suitably chosen separation. At $R=6a_0$ the numerical integral is
\begin{equation}\label{eq:equalz-sigma}
\boxed{S_{sp^3-sp^3}(R=6a_0)=0.71942.}
\end{equation}
For reference, the same-scale transverse $p_x$--$p_x$ pair has $S_{\pi}=0.46800$ at this separation. It would be incorrect to proclaim $S_\sigma>S_\pi$ universally: at much smaller $R$, where the diffuse $sp^3$ functions overlap heavily, cancellations and radial geometry can reverse that particular inequality. These are \emph{overlap} comparisons, not bond strengths or realistic $Z=1$ molecules.
\plot{09_revised_sigma_pi_overlap}{0.97}{Overlap as a function of distance for equal-radial-size $sp^3$--$sp^3$ and $p_x$--$p_x$ pairs (left), and for carbon-sized $sp^2$ $\sigma$ and transverse $\pi$ ingredients (right). A well-defined radial basis matters as much as the angular labels.}

\subsection{A proper $sp^3$--$sp^3$ VB bond in ethane}
Take the C--C bond in ethane, with the \emph{supplied} C--C separation $R=1.543\,\text{\AA}=2.916a_0$. Use $\zeta_C=1.625$ in the inward-facing $sp^3$ hybrids of Eq.~\eqref{eq:sto-sigma}, with $\zeta$ replaced by $\zeta_C$. Their calculated atomic-orbital overlap is
\begin{equation}\label{eq:ethane-overlap}
\boxed{S_{\sigma,\mathrm{C-C}}=0.64659.}
\end{equation}
The corresponding \emph{VB covalent singlet} is $\Phi_{\sigma,\mathrm{cov}}^{S=0}$ of Eq.~\eqref{eq:VBpair-general}, with $a=a_{C_1,sp^3}$ and $b=b_{C_2,sp^3}$. The other six localized C--H pairs and the four carbon $1s$ core electrons are not decorations: their antisymmetrization, electron repulsion, and core exclusion change the quantitative C--C interaction. The full-valence calibration finds one C--C localized occupied pair with populations $50.0\%/49.9\%$ and a C--C bond index
\begin{equation}\boxed{W_{\mathrm{C-C}}(\mathrm{C_2H_6})=1.026,\quad W_{\mathrm{C-H}}\simeq0.991.}\end{equation}
This reproduces the chemically expected \emph{single $\sigma$ pair} without claiming that any individual eH orbital eigenvalue is a bond energy.

For this equal-center VB pair, adding two ionic configurations provides $|C_1^-C_2^+\rangle$ and $|C_1^+C_2^-\rangle$ resonance channels. They have equal energy by molecular symmetry. The spin triplet instead carries the exchange node $\psi_-(\mathbf r,\mathbf r)=0$ and cannot mix with a doubly occupied single-center singlet through a spin-independent hopping operator.
\FloatBarrier

\section{$p$--$p$ $\pi$ bonding: a second spin-coupled pair in ethylene}
\subsection{The same equal-$Z=1$ test, but transverse $p$ orbitals}
Retain the internuclear $z$ axis and the equal nodeless $\zeta=0.5$ valence radial scale, but choose $2p_x$ functions on the two atoms. Then
\begin{equation}\label{eq:sto-pi}
\boxed{a_\pi=\sqrt{\frac{\zeta^5}{\pi}}\,x\e^{-\zeta r_A},\qquad
b_\pi=\sqrt{\frac{\zeta^5}{\pi}}\,x\e^{-\zeta r_B}.}
\end{equation}
Both orbitals have a node on the plane $x=0$, which contains the bond axis. Their overlap is
\begin{equation}\label{eq:sto-pi-overlap}
\begin{aligned}
S_\pi&=\frac{\zeta^5}{\pi}\int x^2\e^{-\zeta(r_A+r_B)}d^3r\\
&=\zeta^5\int_{-\infty}^{\infty}dz\int_0^\infty\rho^3\e^{-\zeta(r_A+r_B)}d\rho
=\boxed{0.46800\quad(R=6a_0,\,\zeta=0.5).}
\end{aligned}
\end{equation}
The $\pi$ pair is \emph{another} singlet of the Heitler--London form, now $\Phi_{\pi,\mathrm{cov}}^{S=0}$ with $a=a_\pi,b=b_\pi$. The two-electron spin algebra is identical to the $\sigma$ case; the difference lies in the nodal structure and in the corresponding overlap, Coulomb exchange, and charge-transfer integrals.
\plot{10_revised_sigma_pi_shapes}{0.96}{Signed atomic $sp^2$ bond-directed and transverse $p$ ingredients in an ethylene-like geometry. A $\pi$ VB pair occupies regions on both sides of a nodal plane through the bond, whereas the $\sigma$ pair concentrates along the internuclear direction.}
\plot{15_vb_sigma_pi_pair_amplitudes}{0.98}{Sections through normalized \emph{two-electron spatial} VB wavefunctions for equal-size $sp^3$ $\sigma$ and $p_x$ $\pi$ ingredients. Columns give the covalent singlet and triplet spatial symmetries. The dashed line is $z_1=z_2$ in the specified fixed transverse slices; the triplet vanishes when the two full electron positions coincide. Color is signed amplitude, not electron density.}

\subsection{Ethylene: one $sp^2$ $\sigma$ pair plus one transverse $p$ $\pi$ pair}
For ethylene at the \emph{input} $R_{\mathrm{C=C}}=1.339\,\text{\AA}=2.530a_0$, a bond-directed $sp^2$--$sp^2$ pair and a transverse $p$--$p$ pair made with the same $\zeta_C=1.625$ give
\begin{equation}\label{eq:ethene-overlap}
\boxed{S_{\sigma,sp^2-sp^2}=0.77052,\qquad S_{\pi,p-p}=0.27039.}
\end{equation}
The full-valence occupied space localizes into four C--H $\sigma$ orbitals and \emph{two distinct C--C orbitals}. The C--C $\pi$ orbital lies completely in the transverse $p$-orbital symmetry block of this planar model. The second C--C orbital is $\sigma$-like. The resulting total C--C bond index is
\begin{equation}\boxed{W_{\mathrm{C=C}}(\mathrm{C_2H_4})=2.022,}\end{equation}
consistent with the expected \emph{two electron pairs} rather than a single two-electron bond whose overlap happens to be larger. The distinction between one $\sigma$ pair and one $\pi$ pair is essential in VB theory.
\plot{17_ethene_sigma_pi_localized_orbitals}{0.96}{The two occupied C--C orbitals from the full-valence model after population-based localization. One resides in the transverse $p_z$ sector (the $\pi$ pair); the other is $\sigma$-like. Their signed AO coefficients are \emph{not} individually unique physical observables or spin-coupled VB amplitudes.}

At the VB level a single dominant Lewis-type structure has the schematic form
\begin{equation}\label{eq:ethene-vb}
\boxed{\Psi_{\mathrm{ethylene}}^{\mathrm{VB}}\sim\mathcal A\left[\Phi_{\mathrm{C-C},\sigma}^{S=0}\,
\Phi_{\mathrm{C-C},\pi}^{S=0}\,
\prod_{\mathrm{four\ C-H}}\Phi_{\mathrm{C-H}}^{S=0}\right],}
\end{equation}
with the C cores included and a suitable many-electron normalization. This product is schematic, not an assertion of energetic additivity. In particular the electron pairs need not be dynamically independent, and other VB structures can mix.

Rotating one $p$ orbital around the C--C axis by torsion angle $\theta$ gives a leading $\pi$ overlap proportional to $\cos\theta$. Accordingly its effective $t_\pi\sim\cos\theta$ and, in a strong-repulsion pair model, its virtual-ionic stabilization approximately follows $\cos^2\theta$. The $\sigma$ pair survives the same torsion. This is a concrete VB explanation for the distinct angular sensitivity of a $\pi$ bond; it is \emph{not} a quantitative rotation barrier prediction.
\FloatBarrier

\section{Polarized $sp^3$--$sp^3$ $\sigma$ bond: C--O in methanol}
\subsection{Unequal orbital sizes and one-electron level mismatch}
An idealized C--O $sp^3$--$sp^3$ pair uses the nodeless valence functions with $\zeta_O/\zeta_C=1.400$ and radial sizes $\zeta_C=1.625$, $\zeta_O=2.275$. At the supplied methanol C--O separation $R=1.427\,\text{\AA}=2.697a_0$, the inward-facing hybrids have
\begin{equation}\label{eq:CO-S}
\boxed{S_{\mathrm{C-O},sp^3-sp^3}=0.54451.}
\end{equation}
In this elementary hybrid basis, $\varepsilon_{C,sp^3}=\tfrac14\varepsilon_{C,2s}+\tfrac34\varepsilon_{C,2p}=-13.90$ eV and $\varepsilon_{O,sp^3}=-19.98$ eV. Oxygen is a deeper atomic \emph{site}; the relative sign of the diagonal mismatch matters for the one-electron polarization. An independent two-hybrid eH projection gives approximately $64\%$ oxygen weight in the lower orthogonalized orbital, but this is not a VB ionic-structure weight.

Including all the valence orbitals changes the quantitative partition. In the full-valence model one localized occupied C--O $\sigma$ orbital has
\begin{equation}\label{eq:CO-ibopop}
\boxed{q_O=0.709,\qquad q_C=0.284,\qquad q_{\rm others}=0.007,\qquad W_{\mathrm{C-O}}=0.837.}
\end{equation}
There are also two strongly O-localized lone-pair orbitals (each $>99\%$ O) and an O--H bond orbital ($\sim69\%$ O). These additional occupied directions are \emph{why a full-valence calculation is important} for interpreting a C--O bond. The population definition is model dependent, and eH overpolarizes some atomic charges; the main robust conclusion is O-biased \emph{covalent} bonding, not $100\%$ ionic transfer.
\plot{11_revised_localized_bonds}{0.96}{Atomic population decompositions of localized occupied orbitals for methane, ethylene, and methanol. The two C--C localized orbitals in ethylene map to $\sigma/\pi$ VB pair regions; methanol shows a polarized C--O pair alongside O lone pairs and other bonds.}

\subsection{Polarization as an explicit VB covalent--ionic resonance problem}
In VB theory, the C--O \emph{singlet} is not just the lower one-electron eigenvector. Use the covalent singlet and two ionic singlets
\begin{equation}
\ket C=\ket{C\!:\!O}_{\rm cov,S=0},\quad
\ket{I_C}=\ket{C^-O^+},\quad
\ket{I_O}=\ket{C^+O^-},\qquad
\ket{\Psi_s}=c_C\ket C+c_{I_C}\ket{I_C}+c_{I_O}\ket{I_O}.
\end{equation}
For atom-centered \emph{overlapping} orbitals these structures are nonorthogonal and obey $\mathsf H_{\rm VB}\mathbf c=E\mathsf S_{\rm VB}\mathbf c$. To expose the physics analytically, first use an \emph{orthonormalized} C/O site basis with site energies $\varepsilon_C,\varepsilon_O$, on-site repulsions $U_C,U_O$, intersite repulsion $V$, and single-electron hopping $t$. Subtract the covalent reference $\varepsilon_C+\varepsilon_O+V$. Ignoring direct Coulomb exchange for this intermediate step gives
\begin{equation}\label{eq:vb-hetero3}
\boxed{H_{S=0}^{\rm VB}=\begin{pmatrix}
0&-\sqrt2\,t&-\sqrt2\,t\\
-\sqrt2\,t&\Delta_C&0\\
-\sqrt2\,t&0&\Delta_O
\end{pmatrix},\qquad
\begin{aligned}
\Delta_C&=U_C-V+\varepsilon_C-\varepsilon_O,\\
\Delta_O&=U_O-V+\varepsilon_O-\varepsilon_C.
\end{aligned}}
\end{equation}
A lower O site energy reduces $\Delta_O$, encouraging the $C^+O^-$ ionic \emph{resonance structure}; a relatively large $U_O$ pushes back against full charge transfer. This competition makes more physical sense than reading polarity from a single $Z_{\rm eff}$ ratio.

When both charge-transfer gaps are positive and $|t|\ll\Delta_{C,O}$, second-order perturbation theory gives
\begin{equation}\label{eq:vb-hetero-second}
\boxed{E_s^{(2)}=-2t^2\left(\frac1{\Delta_C}+\frac1{\Delta_O}\right),\qquad E_t^{(2)}=0\quad\text{(no direct exchange)}.}
\end{equation}
In this orthogonal-site toy, the oxygen electron number in the spin-paired bond is $\langle n_O\rangle=|c_C|^2+2|c_{I_O}|^2$. Its fraction $\langle n_O\rangle/2$ is \emph{not} the same object as the eH localized orbital's $q_O$ in Eq.~\eqref{eq:CO-ibopop}.
\plot{16_vb_polarization_ionic_mixing}{0.96}{Illustrative, exactly diagonalized three-configuration VB singlet with $t/(U-V)=0.18$ and equal $U_C-V=U_O-V$; atomic site mismatch varies along the horizontal axis. Oxygen-like site B becomes the favored ionic configuration as its site energy is lowered. The dashed second-order expression ceases to be controlled near ionic resonance. This is a \emph{dimensionless VB model}, not a fit to methanol.}
\FloatBarrier

\section{Carbon--hydrogen $sp^3$--$1s$ $\sigma$ pairs: methane}
\subsection{An asymmetric orbital pair but four equivalent bonds}
For methane use an inward-directed C $sp^3$ function with $\zeta_C=1.625$ and an H $1s$ function with $\zeta_H=1.300$. At the supplied C--H distance $R=1.087\,\text{\AA}=2.054a_0$,
\begin{equation}\label{eq:ch-new}
\boxed{a_{C,sp^3}=\frac12(s_C+\sqrt3 p_{C,\hat n}),\quad
b_{H,1s}=\left(\frac{\zeta_H^3}{\pi}\right)^{1/2}\e^{-\zeta_Hr_H},\quad S_{C-H}=0.67349.}
\end{equation}
Unlike the original screened-Coulomb calculation, the full-valence model recovers \emph{four equivalent occupied C--H localized bond orbitals} at tetrahedral positions. Each has approximately $50.6\%$ C and $49.4\%$ H population, with $W_{C-H}=1.000$. The carbon population is slightly larger, consistent with modest C$^{\delta-}$--H$^{\delta+}$ polarization, although this method clearly underestimates the size of that preference relative to some correlated population schemes.

The corresponding \emph{VB} picture uses four singlet-paired $C(sp^3)$--$H(1s)$ two-electron functions, antisymmetrized over all eight valence electrons, together with the occupied C core. It is not four unrelated two-electron molecules: the hybrids must be mutually compatible, the electron pairs must obey Pauli exclusion collectively, and the molecular energy cannot be obtained by simply multiplying four isolated pair eigenvalues.

Because the site-energy mismatch between C's inward $sp^3$ average ($-13.90$ eV) and H $1s$ ($-13.60$ eV) is much smaller than in the C--O example, the reference two-hybrid model has a nearly equally shared bonding state. \emph{This is a useful one-electron consistency check, not a quantitative spin-pair energy.} The covalent/ionic VB construction proceeds using the same Eq.~\eqref{eq:vb-hetero3}, now with $C$ and $H$ and a separate specification of $U_C,U_H,V$.
\FloatBarrier

\section{Validation, parameter sensitivity, and what is actually learned}
\subsection{Four chemical case studies versus the discarded models}
\begin{center}\small
\begin{tabular}{p{.27\linewidth}p{.26\linewidth}p{.36\linewidth}}
\toprule
\textbf{Pair or molecule} & \textbf{Numerical check} & \textbf{VB interpretation and boundary}\\\midrule
$Z=1$ nodeless $sp^3$--$sp^3$ $\sigma$ (shape test, $R=6a_0$) & $S_\sigma=0.7194$ & Facing hybrids give a spatially directed, spin-coupled \emph{covalent singlet}; not a bond energy.\\
$Z=1$ $p_x$--$p_x$ $\pi$ (same $R$) & $S_\pi=0.4680$ & Side-on pair with transverse nodal plane; exchange symmetry is the same as the $\sigma$ singlet.\\
Ethane C--C & $W_{CC}=1.026$ & One localized C--C $\sigma$ pair, plus six C--H pairs.\\
Ethylene C=C & $W_{CC}=2.022$, one transverse $p$ orbital & One C--C $\sigma$ singlet \emph{and} one C--C $\pi$ singlet.\\
Methanol C--O & $q_O=0.709$, $W_{CO}=0.837$ & A polarized covalent $\sigma$ pair plus $C^+O^-$/$C^-O^+$ resonances; O lone pairs matter.\\
Methane C--H & $q_C=0.506$, $W_{CH}=1.000$ & Four nearly nonpolar equivalent $sp^3$--$1s$ spin-paired $\sigma$ bonds; polarity remains approximate.\\\bottomrule
\end{tabular}
\end{center}
The input bond lengths were ethane C--C $1.543\,\text{\AA}$, ethylene C=C $1.339\,\text{\AA}$, methanol C--O $1.427\,\text{\AA}$, and methane C--H $1.087\,\text{\AA}$. Their qualitative experimental ordering is used as a \emph{benchmark}, not predicted from minimization of an extended H\"uckel total energy.
\plot{12_revised_bondorder_robustness}{0.95}{Selected full-valence bond indices at $K=1.75$ compared with Lewis bookkeeping (left), and their change across $K$ values (right). The numerical agreement is a \emph{test of one-electron bond topology}, not of the full two-electron VB Hamiltonian.}
\plot{13_revised_polarity_audit}{0.95}{Why the old literal screened-charge model is rejected: its predicted C--O and C--H electron localization can be grossly unphysical. The independent two-hybrid and full-valence calculations are substantially more credible but still model dependent. The published NBO reference bars use a different population definition and are contextual comparisons only.}
\plot{14_revised_polarity_sensitivity}{0.95}{Sensitivity of localized C--O and C--H orbital populations to the Wolfsberg--Helmholtz prefactor. C--O remains oxygen biased; C--H is nearly balanced and less quantitatively reliable. NBO reference values have different definitions and must not be treated as exact targets for L\"owdin populations.}

The robustness sweep $K\in[1.55,2.10]$ gives approximately $W_{CC}(\mathrm{ethane})=1.019$--$1.024$, $W_{CC}(\mathrm{ethylene})=2.009$--$2.023$, and $W_{CO}(\mathrm{methanol})=0.738$--$0.919$. For methanol the localized O fraction of the C--O pair decreases from $0.760$ to $0.661$ across this range, but remains larger than the C fraction. Methane's C fraction varies only from $0.498$ to $0.513$; a weak C--H polarity should not be overinterpreted.

\subsection{A chemically realistic VB calculation requires two-electron integrals}
The shared mathematical object across all the examples is the spin-coupled \emph{two-electron} pair. Given orbitals $a,b$ and a definite Hamiltonian, compute the direct and exchange integrals, covalent--ionic couplings, and ionic diagonal energies, then solve $\mathsf H_{\rm VB}\mathbf c=E\mathsf S_{\rm VB}\mathbf c$. The \emph{sign} of direct Coulomb exchange often favors triplets, whereas virtual ionic configurations often favor singlets; the physical exchange constant comes from the combined Hamiltonian. The actual C--C, C=C, C--O, and C--H \emph{dissociation energies} further require core/valence antisymmetrization, the other bonds and lone pairs, geometry relaxation, and sufficiently correlated total energies.

A quantitative follow-up should use a multiconfigurational or generalized-valence-bond treatment with explicit core exclusion and independently optimized atomic/hybrid orbitals, calibrated against RHF/CASSCF or correlated ab initio calculations. An occupied MO localization alone preserves the one-determinant state and cannot supply those missing pair correlations. The present work therefore provides \emph{validated bond symmetry, population trends, exact toy-model exchange, and a concrete recipe for the true VB calculation}, but does not pretend to predict accurate chemical bond energies.

\highlight{\textbf{Unifying physical picture.} Start with localized atomic \emph{orbitals}, specify a Hamiltonian to compute hopping and Coulomb exchange, impose fermionic antisymmetry to build \emph{spin-coupled covalent pairs}, and add higher-energy \emph{ionic VB structures} as variational or perturbative corrections. The $\sigma/\pi$ distinction is the spatial angular symmetry of the pairs; polarity is their asymmetric covalent--ionic mixing.}

\section{The $1s$ two-center Coulomb integrals}\label{app:1s}
For $\mathrm H_2^+$ use the prolate spheroidal coordinates $\mu=(r_A+r_B)/R$, $\nu=(r_A-r_B)/R$, $\varphi$ around the internuclear axis, with $1\le\mu<\infty$ and $-1\le\nu\le1$,
\begin{equation}\label{eq:prolate}
r_A=\tfrac R2(\mu+\nu),\quad r_B=\tfrac R2(\mu-\nu),\quad
z=\tfrac R2\mu\nu,\quad d^3r=\frac{R^3}{8}(\mu^2-\nu^2)\,d\mu\,d\nu\,d\varphi.
\end{equation}
Because $ab=\pi^{-1}\e^{-R\mu}$,
\begin{align}
S&=\frac{R^3}{4}\int_1^\infty\!\e^{-R\mu}\,d\mu\int_{-1}^1\!(\mu^2-\nu^2)\,d\nu
=\e^{-R}\left(1+R+\tfrac13R^2\right),\\
K_N&=\int ab/r_A\,d^3r
=R^2\int_1^\infty\mu\e^{-R\mu}\,d\mu=(1+R)\e^{-R}.
\end{align}
For $J_N=\int |a|^2/r_B\,d^3r$, use the angular average of the inverse distance from a spherical charge distribution: $\frac{1}{4\pi}\int d\Omega\,|\mathbf r-\mathbf R|^{-1}=1/\max(r,R)$. This gives
\begin{equation}
J_N(R)=\frac4R\int_0^R r^2\e^{-2r}\,dr+4\int_R^\infty r\e^{-2r}\,dr
=\frac{1-(1+R)\e^{-2R}}R.
\end{equation}
Every $1s$ integral is now analytic; the only approximations in the $\mathrm H_2^+$ curves are the \emph{restriction to the two fixed atomic orbitals} and the clamped-nuclei model.

\section{Double-delta derivations and Gaussian integrals}\label{app:delta}
For $V(x)=-g\delta(x-x_A)$, integrate the one-dimensional Schr\"odinger equation across $x_A$ to get $\psi'(x_A^+)-\psi'(x_A^-)=-2mg\psi(x_A)/\hbar^2$. Matching the even and odd double-exponential ansatz at both sites yields $q_\pm=\kappa(1\pm\e^{-q_\pm R})$. At $R>1/\kappa$ the odd positive root is nonzero; below threshold, the only $q_-=0$ solution is non-normalizable. For each decay constant, the identity
\begin{equation}
\int_{-\infty}^\infty\e^{-q|x+R/2|}\e^{-q|x-R/2|}\,dx
=\e^{-qR}(R+q^{-1})
\end{equation}
gives the exact $N_\pm$ in Eq.~\eqref{eq:delta-exact}.

For normalized Gaussian $g_A,g_B$, differentiation gives $g_B''=[(x-R/2)^2/(4\sigma^4)-1/(2\sigma^2)]g_B$. The product identity
\begin{equation}
g_A(x)g_B(x)=\frac{\e^{-R^2/(8\sigma^2)}}{\sqrt{2\pi}\sigma}\e^{-x^2/(2\sigma^2)}
\end{equation}
leads to
\begin{equation}
\mel{g_A}{T}{g_B}=S_G\left(\frac1{8\sigma^2}-\frac{R^2}{32\sigma^4}\right),\quad
\mel{g_A}{-g\delta(x\pm R/2)}{g_B}=-\frac g{\sqrt{2\pi}\sigma}\e^{-R^2/(4\sigma^2)}.
\end{equation}
This proves Eq.~\eqref{eq:gauss-mat} without numerical quadrature. The variational minimum of an even or odd superposition remains an upper bound on the exact \emph{lowest state of that parity}, provided such a discrete state exists.

\section{Historical hydrogenic-hybrid quadrature (superseded shape model)}\label{app:prolate}
The ingredients of all four worked 3D examples are normalized \emph{one-electron} orbitals of the form
\begin{equation}
a=P_A(r_A,z_A,\rho,\varphi)\e^{-\lambda_A r_A},\quad
b=P_B(r_B,z_B,\rho,\varphi)\e^{-\lambda_B r_B},\quad
\lambda_{2s,2p}=Z/2,\quad\lambda_{1s}=Z.
\end{equation}
Here $P$ is the full normalizing polynomial including $C_Z$ (and, for $p_x$, $\rho\cos\varphi$). With Eq.~\eqref{eq:prolate},
\begin{equation}\label{eq:general-prolate}
\begin{split}
\int abF\,d^3r={\frac{R^3}{8}}\int_1^\infty d\mu\int_{-1}^1d\nu\int_0^{2\pi}d\varphi\,&(\mu^2-\nu^2)P_AP_B\,F\\[-2pt]
&\times\exp\left[-\tfrac R2\big((\lambda_A+\lambda_B)\mu+(\lambda_A-\lambda_B)\nu\big)\right],
\end{split}
\end{equation}
where $F=1$, $1/r_A$, or $1/r_B$. For a $\sigma$--$\sigma$ pair the integrand is independent of $\varphi$, giving a factor $2\pi$. For a $p_x$--$p_x$ pair, $P_AP_B\propto\rho^2\cos^2\varphi$ and the azimuthal integral is $\pi$, i.e., half the $\sigma$ factor. The diagonal $|a|^2$ and $|b|^2$ integrals use the same formula with $a=b$ on the appropriate center.

For numerical accuracy, the plotted calculations transform $\mu=1+u/c$ with $c=(\lambda_A+\lambda_B)R/2$, factor out $\e^{-c}$, and use 90-point Gauss--Laguerre quadrature in $u\in[0,\infty)$ together with 110-point Gauss--Legendre quadrature in $\nu\in[-1,1]$. This integrates the full radial tails, avoids a finite Cartesian box, and handles the nuclear $1/r$ factors by retaining the prolate Jacobian. Using the exact atomic eigenvalue identities, the independent evaluations
\begin{equation}
E_BS-Z_AK_A\stackrel{!}{=}E_AS-Z_BK_B
\end{equation}
agree to numerical precision in all tabulated cases. Norm checks give $\braket aa=\braket bb=1$ to the precision quoted.

\begin{table}[htbp]\centering\small
\caption{Illustrative projected one-electron results (Hartree energies, $a_0$ separations). These are mathematical toy-model values, not quantitative bond predictions.}
\begin{tabular}{lrrrrr}
\toprule
Orbitals & $R$ & $S$ & $H_{AB}$ & $t_{\rm orth}$ & $w_B^{(-)}$\\\midrule
$sp^3$--$sp^3$, $Z=1$ & 5.0 & $-0.095004$ & $+0.003010$ & 0.022968 & 0.5000\\
$p_x$--$p_x$, $Z=1$ & 5.0 & $+0.578015$ & $-0.186829$ & 0.034638 & 0.5000\\
C--O $sp^3$--$sp^3$ & 2.7 & $+0.032043$ & $-0.202976$ & 0.102116 & 0.9880\\
C--H $sp^3$--$1s$ & 2.1 & $+0.109688$ & $-0.338150$ & 0.138534 & 0.8573\\
\bottomrule
\end{tabular}
\end{table}

\section{Two-center Slater overlap and ionic VB matrix elements}\label{app:sto}
For the nodeless, orthonormal-on-center STO valence basis used above, write
\begin{equation}
 s_{\zeta}=\sqrt{\frac{\zeta^5}{3\pi}}r e^{-\zeta r},\qquad
 p_{j,\zeta}=\sqrt{\frac{\zeta^5}{\pi}}r_j e^{-\zeta r}.
\end{equation}
For bond-directed hybrids $h_A=\frac12(s_A+\sqrt3\,p_{z,A})$ and
$h_B=\frac12(s_B-\sqrt3\,p_{z,B})$, the overlap is the signed combination
\begin{equation}
 S_\sigma=\frac14\bigl(S_{ss}-\sqrt3 S_{s p_B}+\sqrt3 S_{p_A s}-3S_{p_A p_B}\bigr),
 \qquad S_{\pi}=\langle p_{x,A}|p_{x,B}\rangle.
\end{equation}
These signs are tied to the local $p$ phase and the direction each hybrid points; reversing a basis phase changes overlap signs but not observables. The published numerical quadrature evaluates the corresponding three-dimensional integrals at the stated nuclear geometries. The localized-orbital populations $q_A$ are \emph{not} squared coefficients of a covalent--ionic VB expansion.

For nonorthogonal, normalized atomic orbitals $a,b$ with overlap $S$, a compact singlet basis is
\begin{align}
 |C\rangle&=\frac{|ab\rangle+|ba\rangle}{\sqrt{2(1+S^2)}}\,\chi_s,
 &|I_A\rangle&=|aa\rangle\chi_s,&|I_B\rangle&=|bb\rangle\chi_s.
\end{align}
Its overlap matrix is
\begin{equation}
 \mathsf S_{\rm VB}=\begin{pmatrix}
 1&\sqrt2 S/\sqrt{1+S^2}&\sqrt2 S/\sqrt{1+S^2}\\
 \sqrt2 S/\sqrt{1+S^2}&1&S^2\\
 \sqrt2 S/\sqrt{1+S^2}&S^2&1
 \end{pmatrix}.
\end{equation}
For a specified spin-free two-electron Hamiltonian, the variational energy and structure coefficients follow from $\mathsf H_{\rm VB}c=E\mathsf S_{\rm VB}c$. This is a three-configuration \emph{variational} VB treatment. At large separation, eliminating higher-energy ionic configurations perturbatively recovers a superexchange-like singlet stabilization, but the numerical $U$ and $t$ cannot be obtained from the one-electron bond index alone.

\section{Two-electron direct and exchange integrals for arbitrary hybrids}\label{app:twoelectron}
For any two normalized real orbitals $a,b$ with $S=\langle a|b\rangle$, and a symmetric two-electron Hamiltonian $H=h(1)+h(2)+r_{12}^{-1}+V_{NN}$,
\begin{align}
D&=H_{aa}+H_{bb}+J_{ee}+V_{NN},\\
X&=S(H_{ab}+H_{ba})+K_{ee}+V_{NN}S^2,\\
J_{ee}&=\int d^3r_1\int d^3r_2\frac{|a(\mathbf r_1)|^2|b(\mathbf r_2)|^2}{|\mathbf r_1-\mathbf r_2|},\\
K_{ee}&=\int d^3r_1\int d^3r_2\frac{a(\mathbf r_1)b(\mathbf r_1)a(\mathbf r_2)b(\mathbf r_2)}{|\mathbf r_1-\mathbf r_2|}.
\end{align}
For the $sp^3$ $\sigma$ case, substitute $a=a_\sigma$ and $b=b_\sigma$ from Eq.~\eqref{eq:sto-sigma}. For the $p_x$ $\pi$ case, the product at each electron coordinate is
\begin{equation}
a_\pi(\mathbf r)b_\pi(\mathbf r)
=C_1^2x^2\exp[-(r_A+r_B)/2],
\end{equation}
and hence $K_{ee,\pi}$ has the explicit integrand
\begin{equation}
K_{ee,\pi}=C_1^4\iint \frac{x_1^2x_2^2\exp[-(r_{1A}+r_{1B}+r_{2A}+r_{2B})/2]}{r_{12}}\,d^3r_1d^3r_2.
\end{equation}
The two-electron kernels are \emph{not} evaluated by the one-electron $H_{ab}=E_BS-Z_AK_A$ formula. For numerical work, the direct $J_{ee}$ can be evaluated as the electrostatic interaction of two orbital densities; the exchange $K_{ee}$ can be evaluated as the Coulomb self-energy of the \emph{transition density} $\rho_{ab}=ab$, e.g., by solving Poisson's equation for this signed source and integrating $\int\rho_{ab}\Phi_{ab}$. Gaussian integral expansion or elliptical-coordinate methods provide other routes. In this edition we report \emph{one-electron} numerical $\sigma/\pi$ projections and give the complete Heitler--London \emph{two-electron} expressions without inventing uncomputed binding curves.

\section{Reproducibility, conventions, and reading}
The figures are generated by the companion \texttt{make\_figures.py}, \texttt{make\_revised\_figures.py}, and \texttt{make\_vb\_figures.py} scripts. The $\mathrm H_2^+$ example uses analytic fixed-$1s$ integrals; double-delta finite-$R$ roots are found by bracketed numerical root solving; Gaussian widths are variationally optimized; hydrogenic $2s/2p$ two-center integrals use deterministic Gauss--Laguerre and Gauss--Legendre quadrature. Numeric outputs are recorded in \texttt{numerical\_results.json} and \texttt{physical\_revision\_results.json}. Only the \emph{one-electron} matrix elements and orbital populations are numerically evaluated for the higher-orbital examples. The replacement higher-orbital analysis uses spatial STO overlaps and a calibrated semiempirical full-valence Hamiltonian, not screened point-charge Coulomb projections. It does not include self-consistent two-electron correlation, and the VB Coulomb integrals are specified but not numerically solved for polyatomic bonds.

For background and conventions: Heitler and London, \emph{Wechselwirkung neutraler Atome und hom\"oopolar Bindung nach der Quantenmechanik}, Z. Phys. \textbf{44}, 455--472 (1927); A. Szabo and N. Ostlund, \emph{Modern Quantum Chemistry} (Dover); P. W. Atkins and R. Friedman, \emph{Molecular Quantum Mechanics}; J. C. Slater, \emph{Atomic Shielding Constants}, Phys. Rev. \textbf{36}, 57--64 (1930). Hydrogenic orbital normalization conventions are cross-checked against \href{https://chem.libretexts.org/Bookshelves/Physical_and_Theoretical_Chemistry_Textbook_Maps/Supplemental_Modules_(Physical_and_Theoretical_Chemistry)/Quantum_Mechanics/09._The_Hydrogen_Atom/Chapter_7.__The_Hydrogen_Atom}{Chemistry LibreTexts, Hydrogen Atom: Orbital Formulas}; the Slater-rule screened-charge arithmetic is a pedagogical estimate, not a modern atomic calculation.
\end{document}
